\fontsize{13pt}{22.5pt}\selectfont

\section{KẾT LUẬN}

\subsection{Kết quả đạt được}
\indent Đồ án đã hoàn thành các mục tiêu đặt ra ban đầu. Kết quả chính là ba môi trường
huấn luyện game 3D mang tính chất khác nhau trên nền tảng Unity ML-Agents. Football Table
lo phần đối kháng một đối một, hành động liên tục, tự chơi. Capture The Flag lo phần hợp
tác hai tác nhân với phần thưởng nhóm và chương trình học theo bậc. Cross The Road lo phần
điều hướng đơn tác nhân với phần thưởng thưa. Với từng môi trường, đề tài thuật lại trọn
vẹn quá trình thiết kế không gian quan sát, không gian hành động và hàm phần thưởng, kể cả
những phương án đã thử rồi phải bỏ.

\indent Phần tri thức thiết kế thu được từ quá trình đó có giá trị không kém bản thân ba
môi trường. Capture The Flag là trường hợp điển hình. Đặc tả ban đầu của nó trừng phạt
đúng cái trạng thái phối hợp mà nhiệm vụ đòi hỏi, tức hai tác nhân đứng tách nhau ở hai phía
cửa, nên tác nhân mắc kẹt ở đáy phần thưởng nhóm bất kể chạy thuật toán nào. Lỗi đó không
hiện ra thành lỗi chương trình, nó chỉ hiện thành một đường cong phẳng, và chỉ lộ diện khi
ta ngồi xem hành vi thật. Bài học rút ra ở đây là phải kiểm tra xem hàm phần thưởng có vô tình
phạt chính hành vi đích hay không, và nó áp dụng được cho bất kỳ môi trường hợp tác nào, chẳng
riêng gì môi trường này.

\indent Ba môi trường đã được đóng thành một gói Unity Package Manager, kèm $15$ tệp cấu hình
huấn luyện, bộ công cụ Editor và trình hướng dẫn cài đặt bảy bước. Nhờ vậy chúng
tồn tại như một sản phẩm dùng lại được, chứ không chỉ là mã nguồn của một đồ án. Đề tài
cũng bổ sung vào ML-Agents hai thuật toán vốn không có sẵn là MAPPO và DQN dưới dạng trainer
mở rộng, đồng thời hiện thực hóa quy trình nối môi trường Unity với một thư viện
học tăng cường bên ngoài qua Python Low-Level API.

\indent Cả $15$ lần huấn luyện đã xác nhận cả ba môi trường đạt yêu cầu theo ba tiêu chí
đặt ra ở Chương~\ref{ch:results}. Chúng \textit{học được}: ở cả ba môi trường đều có ít
nhất bốn trong năm thuật toán tiến bộ rõ rệt so với lúc khởi tạo. Chúng \textit{phân biệt
được}: trên Cross The Road, phần thưởng cuối trải từ $0{,}846$ của SAC xuống $0{,}336$ của
DQN; còn trên Capture The Flag, bốn thuật toán về đích trong dải $0{,}800$--$0{,}840$
trong khi SAC sụp đổ ở $0{,}183$. Thứ hạng ở cả hai nơi đều giải thích được bằng cơ chế của
bài toán. Chúng cũng \textit{không có lối tắt}: không lần chạy nào đạt phần thưởng
sát trần mà lại không giải được nhiệm vụ khi ta quan sát trực tiếp trong ứng dụng trình
diễn.

\indent Riêng Football Table bộc lộ một giới hạn cần thừa nhận thẳng thắn. Cả bốn thuật
toán ghép được với tự chơi đều tăng ELO từ $36$ tới $64$ điểm so với phiên bản cũ của chính
mình, cho thấy cơ chế tự chơi hoạt động đúng. Nhưng biên độ ấy quá hẹp để môi
trường phân định được các thuật toán với nhau. Ở cấu hình hiện tại, Football Table chỉ đạt
hai trong ba tiêu chí. Muốn dùng nó làm thước đo so sánh thì cần ngân sách bước lớn hơn,
hoặc một đối thủ mốc cố định thay cho tự chơi thuần túy.

\indent Có một quan sát xuyên suốt cả ba môi trường ủng hộ đúng luận điểm trung tâm của đề
tài: chênh lệch của cùng một thuật toán giữa hai môi trường luôn lớn hơn chênh lệch giữa
các thuật toán trong cùng một môi trường. Khoảng cách giữa thành công của SAC trên Cross
The Road và sự sụp đổ của chính nó trên Capture The Flag lớn hơn mọi khoảng cách nội bộ.
Nói cách khác, đặc tính bài toán, thứ do thiết kế môi trường quyết định, chi phối kết quả mạnh
hơn lựa chọn thuật toán. Và đó chính là lý do khâu xây dựng môi trường đáng được
đầu tư nhiều hơn mức thường thấy.

\indent Bên cạnh phần thực nghiệm so sánh, đề tài còn dựng hoàn chỉnh một ứng dụng trình
diễn tương tác, trong đó cả $15$ mô hình đã huấn luyện được triển khai thẳng vào game
dưới định dạng ONNX. Người dùng ở đó chọn môi trường, chọn một trong ba chế độ chơi (tự điều
khiển, tác nhân tự chơi, chơi cùng tác nhân), chọn mô hình cụ thể cho tác nhân, và
tra cứu ngay đường cong huấn luyện của cả năm thuật toán trên giao diện menu. Danh mục mô
hình cùng dữ liệu biểu đồ đều sinh ra từ một kịch bản duy nhất chạy trên thư mục kết quả,
nên phần trình diễn và phần số liệu của Chương~\ref{ch:results} không thể lệch nhau. Các
luồng chức năng đã được kiểm thử trong Unity Editor và đều chạy đúng
(Mục~\ref{sub:app-testing}). Kết quả này chứng minh việc nhúng mô hình học tăng cường vào
một sản phẩm game hoàn chỉnh là khả thi, với chi phí suy luận không đáng kể.

\subsection{Đóng góp của đề tài}
\indent Nếu phải xếp các đóng góp theo mức hữu ích, tôi đặt bộ môi trường lên đầu. Ba môi
trường ba chiều phủ ba dạng bài toán học tăng cường phổ biến, đóng thành một gói Unity
Package Manager cài được vào dự án bất kỳ, kèm $15$ cấu hình huấn luyện đã chạy thật,
bộ công cụ Editor và một trình hướng dẫn cài đặt bảy bước biết tự kiểm tra trạng thái.
Điều khiến nó khác việc chỉ công bố mã nguồn là ba môi trường này đã qua kiểm chứng: mỗi
môi trường đã chịu năm thuật toán khác hẳn nhau về nguyên lý, nên người dùng lại biết
trước rằng nó học được, nó phân biệt được, và nó không có lối tắt, những điều mà một môi
trường mới dựng chẳng bao giờ bảo đảm nổi. Chính đồ án này tiêu thụ gói đó qua ranh giới
assembly y như một dự án bên ngoài, nên phần trình diễn hóa thành phép kiểm thử thường
trực cho gói.

\indent Kế đến là phần tri thức thiết kế môi trường, thứ mà tài liệu về học tăng cường
thường lướt qua. Đề tài ghi lại đầy đủ quá trình đặc tả cho ba dạng bài toán khác nhau,
gồm cả những phương án thất bại và dấu hiệu nhận ra chúng: một hàm phần thưởng trừng phạt
chính hành vi đích; một tác nhân thiếu quan sát về đồng đội nên không định được thời điểm
phối hợp; một khoản phạt thời gian chia cho không khi giới hạn bước bằng không. Đi cùng là
các quyết định vận hành ảnh hưởng lớn nhưng ít được bàn (nhân khu vực huấn luyện để tăng thông
lượng, dựng biến thể không gian hành động rời rạc cho thuật toán dựa trên giá trị), và hai con
đường đưa một thuật toán không có sẵn vào Unity: gói trainer mở rộng đăng ký
MAPPO cùng DQN vào chính chương trình huấn luyện của ML-Agents, và lớp bọc môi trường qua
\texttt{mlagents\_envs} để nối với Stable-Baselines3.

\indent Kết quả so sánh năm thuật toán tuy chỉ là sản phẩm phụ của khâu kiểm chứng nhưng
vẫn tham khảo trực tiếp được cho người làm game. Đáng chú ý nhất là một kết quả phủ định:
trên bài toán hợp tác có tác nhân đồng nhất, PPO thường bám sát cả MAPPO lẫn MA-POCA trong
dải chênh lệch chưa tới một phần trăm, nghĩa là chi phí tích hợp một thuật toán đa tác
nhân chuyên biệt chưa chắc được đền đáp. Đi kèm nó là một kết quả gây ngạc nhiên theo
chiều ngược lại: MAPPO vượt PPO một khoảng lớn trên chính môi trường đơn tác nhân, nơi lý
thuyết bảo rằng phê bình tập trung của nó lẽ ra vô hiệu.

\indent Đóng góp cuối cùng là ứng dụng trình diễn với ba chế độ chơi, cơ chế chọn mô hình
cho phép đặt cạnh nhau cả $15$ chính sách đã huấn luyện, và hai tab số liệu huấn luyện
nhúng thẳng trong game. Giá trị của nó không nằm ở phần game, mà ở chỗ nó là công cụ soi
lỗi đặc tả. Quan sát hành vi thật là cách duy nhất phát hiện những sai sót mà đường cong
phần thưởng giấu đi, và cũng là cách khép kín quãng đường từ một tệp cấu hình YAML tới một
hành vi mà người chơi cảm nhận được bằng tay: chính bảng xếp hạng ở Chương~\ref{ch:results}
có thể kiểm chứng lại bằng mắt, bằng cách cho hai mô hình khác thuật toán chạy trên cùng
một nhiệm vụ. Toàn bộ mã nguồn môi trường, cấu hình huấn luyện và dữ liệu kết quả đều được
công bố dưới dạng mở.

\subsection{Hạn chế}
\label{sub:limits}
\indent Hạn chế nặng nhất của đề tài: mỗi tổ hợp môi trường và thuật toán chỉ chạy với đúng
một seed. Nghĩa là mọi chênh lệch báo cáo ở Chương~\ref{ch:results} đều chưa tách được
phần do bản chất thuật toán khỏi phần do may rủi của một lần khởi tạo. Những khoảng cách lớn
(DQN xếp cuối trên Cross The Road, SAC sụp đổ trên Capture The Flag) nhiều khả năng vẫn đứng
vững; nhưng các thứ hạng sát nhau thì đừng đọc như một kết luận thống kê. Đây là
hệ quả của quỹ thời gian tính toán chứ không phải của một lựa chọn phương pháp luận, và
nếu chỉ được sửa đúng một điều trong đồ án này thì tôi sẽ sửa điều đó.

\indent Hạn chế thứ hai nằm ở độ phân giải của phép đo. Ở Football Table, bốn thuật toán
tự chơi kết thúc trong dải ELO $1237$--$1260$. Ở Capture The Flag, bốn thuật toán dẫn đầu
nằm trong dải phần thưởng $0{,}800$--$0{,}840$. Cả hai khoảng cách đều nhỏ hơn dao động
của chính các đường cong, nên ở cấu hình hiện tại hai môi trường này không phân biệt nổi
các thuật toán trong nhóm dẫn đầu. Chỉ Cross The Road cho một phổ đủ rộng để việc xếp hạng
có nghĩa. Muốn hai môi trường kia nói được nhiều hơn thì phải làm chúng khó lên, chứ không
phải chạy chúng lâu hơn.

\indent Còn một ô của lưới thực nghiệm nằm ngoài phép so sánh: DQN trên Football Table
chạy không có tự chơi, vì ML-Agents báo lỗi khi ghép trainer off-policy với bộ điều phối
đối kháng. Lần chạy đó vì thế không có ELO và học trong điều kiện khác hẳn bốn thuật toán
kia, nên mọi con số của nó ở môi trường này chỉ đọc được như một tham chiếu rời.

\indent Sau cùng, đề tài không giải thích được hai hiện tượng mà chính số liệu của mình
phơi ra. Hiện tượng thứ nhất là MAPPO vượt PPO trên môi trường đơn tác nhân, nơi hai thuật
toán chỉ khác nhau ở phần phê bình tập trung, thứ lẽ ra vô hiệu khi nhóm có mỗi một thành
viên. Hiện tượng thứ hai là kiểu thất bại của SAC trên Capture The Flag: entropy sụp xuống
gần không rồi bật ngược lên đúng mức cực đại và nằm đó tới cuối. Cả hai đều cần thêm thực
nghiệm mới kết luận được; mà với một hạt giống cho mỗi ô thì cũng chưa loại được khả năng
hiện tượng thứ nhất chỉ là dao động ngẫu nhiên.

\subsection{Hướng phát triển}
\indent Việc phải làm trước tiên là chạy lại toàn bộ $15$ tổ hợp với nhiều hạt giống,
bởi chừng nào chưa có phương sai giữa các lần chạy thì mọi so sánh trong đồ án vẫn treo lơ
lửng. Song song, cần làm Football Table và Capture The Flag khó lên, vì ở cấu hình hiện
tại chúng thiếu độ phân giải để tách các thuật toán dẫn đầu. Với Capture The Flag, hướng
cụ thể là thêm bài trong chương trình học hoặc buộc hai tác nhân phối hợp chặt hơn, chứ
không phải kéo dài thời lượng huấn luyện.

\indent Ngay sau đó là hai câu hỏi mà đề tài đặt ra nhưng chưa trả lời nổi. Thứ nhất, vì
sao MAPPO vượt PPO trên môi trường đơn tác nhân? Một thực nghiệm tách riêng phần phê bình
tập trung, chạy trên nhiều hạt giống, sẽ cho biết đây là hiệu ứng thật hay chỉ là may rủi
của một lần khởi tạo. Thứ hai, cơ chế nào đứng sau kiểu thất bại của SAC trên bài toán hợp
tác? Giả thuyết khớp với dữ liệu là bộ đệm phát lại lưu lẫn lộn kinh nghiệm của hai bài
học trong chương trình học, và cách kiểm chứng trực tiếp là chạy lại với bộ đệm tách riêng
theo từng bài. Ranh giới giữa chia sẻ tham số và attention cũng đáng khảo sát, bằng một
môi trường hợp tác có hai tác nhân vai trò bất đối xứng, hoặc có tác nhân rời nhóm giữa
lượt chơi.

\indent Ở tầm xa hơn, bộ so sánh có thể mở rộng bằng những thuật toán như DreamerV3 hay QMIX,
việc bổ sung giờ đã rẻ hơn nhiều nhờ cơ chế trainer mở rộng đã dựng sẵn; còn ba môi trường thì
có thể nâng độ phức tạp lên gần với game thương mại. Riêng ứng dụng trình
diễn cần một bản chạy độc lập, WebGL hoặc desktop, kèm bảng xếp hạng người chơi đấu với
tác nhân và một khảo sát trải nghiệm người dùng. Đó là cách duy nhất trả lời được câu hỏi
mà không chỉ số huấn luyện nào chạm tới: hành vi của tác nhân có thực sự ``tự nhiên''
trong mắt người chơi hay không.
